% \iffalse meta-comment
%
%% Copyright (C) 2025 by Marcel Krueger
%%
%% This file may be distributed and/or modified under the
%% conditions of the LaTeX Project Public License, either
%% version 1.3c of this license or (at your option) any later
%% version. The latest version of this license is in:
%%
%% http://www.latex-project.org/lppl.txt
%%
%% and version 1.3 or later is part of all distributions of
%% LaTeX version 2005/12/01 or later.
%
%<*gobble>
\expandafter\ifx\csname @currname\endcsname\empty
\csname fi\endcsname
%</gobble>
%<*driver>
\documentclass[full,a4paper]{l3doc}
\usepackage{booktabs}
\usepackage{lua-unicode-math}
\MakeShortVerb{\|}
\begin{document}
\DocInput{lua-unicode-math.dtx}
% \PrintIndex
% \PrintChanges
\end{document}
%</driver>
%<*gobble>
\fi
%</gobble>
% \fi
%
% \GetFileInfo{lua-unicode-math.dtx}
% \title{The \pkg{lua-unicode-math} package\thanks{This document
%        corresponds to \pkg{lua-unicode-math}~v0.2, dated~2025-12-11.}}
% \author{Marcel Kr\"uger \\ \href{mailto:tex@2krueger.de}{tex@2krueger.de} \\ \url{https://github.com/zauguin/lua-unicode-math}}
% 
% \let\fntname\texttt
% \maketitle
% \begin{documentation}
% Modern fonts are usually provided in OpenType format and are designed for Unicode
% based input. For mathematical fonts this usually means the use of fonts with an OpenType MATH table:
% Fonts containing special metadata needed to make them usable in a mathematical context.
%
% In Lua\TeX{} such fonts have traditionally been loaded with the \pkg{unicode-math} package.
% While this works, is very flexible and allows to use the same document in Xe\TeX{} and Lua\TeX{} it has performance issues and it sometimes has unexpected interactions with the use of math versions. The \pkg{lua-unicode-math} is a specific Lua\LaTeX{} specific alternative which aims for higher performance and better integration with native Lua\TeX{} features.
%
% \section{Usage instructions}
% \subsection{Font packages}
% For most Opentype the recommended way to load them with \pkg{lua-unicode-math} is to use a dedicated package.
% Currently the following packages are shipped with \pkg{lua-unicode-math}:
% \begin{tabular}{ll}
% \toprule
% Font & Package \\
% \midrule
% Latin Modern Math & lum-lmodern \\
% New Computer Modern Math & lum-newcomputermodern \\
% STIX2 & lum-stix2 \\
% XITS & lum-xits \\
% TeX Gyre Pagella Math & lum-pagella \\
% TeX Gyre DejaVu Math & lum-dejavu \\
% TeX Gyre Bonum Math & lum-bonum \\
% TeX Gyre Schola Math & lum-schola \\
% TeX Gyre Termes Math & lum-termes \\
% Fira Math & lum-fira \\
% GFS Neohellenic Math & lum-gfsneohellenic \\
% Erewhon Math & lum-erewhon \\
% XCharter Math & lum-xcharter \\
% Concrete Math & lum-concrete \\
% \bottomrule
% \end{tabular}
%
% \subsection{Loading fonts by name}
% If you want to use a custom font, you can load \pkg{fontspec} and \pkg{lua-unicode-math}
% using
% \begin{verbatim}
% \usepackage{fontspec,lua-unicode-math}
% \end{verbatim}
% This will load \fntname{Latin Modern Math} by default. Another math font can be loaded using \cs{setmathfont} using the same options as \pkg{fontspec}'s \cs{newfontfamily}.
% For example, you can use to to configure the current math font using
% \begin{verbatim}
% \setmathfont[AutoFakeBold=1]{Latin Modern Math}
% \end{verbatim}
%
% \subsection{Writing maths}
% There are two ways of entering math: You can directly input Unicode math symbols or use regular \LaTeX{} commands for symbols.
% All Unicode symbols are supported with the same commands as in \pkg{unicode-math}. For a full list see \texttt{texdoc unimath-symbols}.
% \end{documentation}
% \begin{implementation}
% \section{Implementation}
%    \begin{macrocode}
\ProvidesExplPackage
  {lua-unicode-math}
  {2025-12-11}
  {0.2}
  {Opentype Math support for LuaLaTeX}

%<@@=l_uni_math>
\int_new:N \g__l_uni_math_font_count_int
\tl_new:N \l__l_uni_math_main_family_tl
\tl_new:N \l__l_uni_math_script_family_tl
\tl_new:N \l__l_uni_math_scriptscript_family_tl

\cs_generate_variant:Nn \tl_if_eq:nnT {o}

\msg_new:nnn { l_uni_math } { unicode-math-suppressed } {
  You~tried~to~load~both~lua-unicode-math~and~unicode-math~
  in~the~same~document.~This~is~not~supported,~unicode-math~
  will~be~suppressed.~There~is~a~good~chance~that~this~will~
  break~your~document.~Change~your~document~to~only~use~lua-unicode-math~
  so~solve~this.
}
\msg_new:nnn { l_uni_math } { unicode-math-loaded } {
  You~tried~to~load~lua-unicode-math~while~unicode-math~
  was~already~loaded.~This~does~not~work.~Please~avoid~loading~
  unicode-math.~If~that~is~not~possible~and~you~are~feeling~adventurous~
  you~can~try~loading~the~lua-unicode-math~package~at~the~beginning~
  of~your~document~instead~to~suppress~unicode-math.
}
\disable@package@load{unicode-math} {
  \msg_warning:nn { l_uni_math } {unicode-math-incompatible }
}
\IfPackageLoadedTF {unicode-math} {
  \msg_critical:nn { l_uni_math } {unicode-math-loaded }
} {}


\cs_if_exist:NF \DeclareMathScriptfontMapping {
  \cs_new:Npn \DeclareMathScriptfontMapping #1 #2 #3 #4 #5 #6 {
    \cs_set:cpx { scriptfont@@ #1 / #2 } { #3 / #4 }
    \cs_set:cpx { scriptscriptfont@@ #1 / #2 } { #5 / #6 }
  }
  \cs_set:Npn \__l_uni_math_family_variant:nw #1 #2/#3/ {
    \ifcsname #1@@#2/#3 \endcsname
      \lastnamedcs
    \else
      #2/#3
    \fi
    /
  }

% Hook into the kernel to allow different families in scriptfonts
  \cs_set:Npn \getanddefine@fonts #1 #2 {
    \xdef\font@name{\csname \string#2/\tf@size\endcsname}
    \pickup@font \let\textfont@name\font@name
    \xdef\font@name{\csname \exp_last_unbraced:Nno \__l_uni_math_family_variant:nw {scriptfont} {\string#2/\sf@size}\endcsname}
    \pickup@font \let\scriptfont@name\font@name
    \xdef\font@name{\csname \exp_last_unbraced:Nno \__l_uni_math_family_variant:nw {scriptscriptfont} {\string#2/\ssf@size}\endcsname}
    \pickup@font
    \edef\math@fonts{\math@fonts
                 \textfont#1\textfont@name
                 \scriptfont#1\scriptfont@name
                 \scriptscriptfont#1\font@name}
  }
}

\hook_gput_code:nnn { package/fontspec/after } {.} {
  \bool_gset_false:N \g__fontspec_math_bool

  \NewDocumentCommand \setmathfont { O{} m O{} } {
    \int_incr:N \g__l_uni_math_font_count_int
    \exp_args:Nc \newfontfamily
      { g__l_uni_math_font_ \int_use:N \g__l_uni_math_font_count_int _text_font }
      { #2 }
      [ #1, #3, Script = Math, Renderer = Base ]
    \tl_set_eq:NN \l__l_uni_math_main_family_tl \l_fontspec_family_tl

    \exp_args:Nc \newfontfamily
      { g__l_uni_math_font_ \int_use:N \g__l_uni_math_font_count_int  _script_font }
      { #2 }
      [ #1, #3, Script = Math, Renderer = Base, Style = MathScript ]
    \tl_set_eq:NN \l__l_uni_math_script_family_tl \l_fontspec_family_tl

    \exp_args:Nc \newfontfamily
      { g__l_uni_math_font_ \int_use:N \g__l_uni_math_font_count_int  _scriptscript_font }
      { #2 }
      [ #1, #3, Script = Math, Renderer = Base, Style = MathScriptScript ]
    \tl_set_eq:NN \l__l_uni_math_scriptscript_family_tl \l_fontspec_family_tl

    \DeclareMathScriptfontMapping {TU} {\l__l_uni_math_main_family_tl} {TU} {\l__l_uni_math_script_family_tl} {TU} {\l__l_uni_math_scriptscript_family_tl}

    \exp_args:NnnV \DeclareSymbolFont {lummain} {TU} \l__l_uni_math_main_family_tl {m} {n}
    \exp_args:NnnnV \SetSymbolFont {lummain} {bold} {TU} \l__l_uni_math_main_family_tl {b} {n}
  }

  \cs_set:Nn \__fontspec_setmainfont_hook:nn
    {
      \tl_if_eq:onT {\g__fontspec_mathrm_tl} {\rmdefault}
        {
          \fontspec_gset_family:Nnn \g__fontspec_mathrm_tl {Renderer=Basic,#1} {#2}
          \__fontspec_setmathrm_hook:nn {#1} {#2}
        }
    }
  \cs_set:Nn \__fontspec_setsansfont_hook:nn
    {
      \tl_if_eq:onT {\g__fontspec_mathsf_tl} {\sfdefault}
        {
          \fontspec_gset_family:Nnn \g__fontspec_mathsf_tl {Renderer=Basic,#1} {#2}
          \__fontspec_setmathsf_hook:nn {#1} {#2}
        }
    }
  \cs_set:Nn \__fontspec_setmonofont_hook:nn
    {
      \tl_if_eq:onT {\g__fontspec_mathtt_tl} {\ttdefault}
        {
          \fontspec_gset_family:Nnn \g__fontspec_mathtt_tl {Renderer=Basic,#1} {#2}
          \__fontspec_setmathtt_hook:nn {#1} {#2}
        }
    }
  \cs_set:Nn \__fontspec_setmathrm_hook:nn
    {
      \SetMathAlphabet \mathrm { normal } \g_fontspec_encoding_tl \g__fontspec_mathrm_tl { \mdseries@rm } \shapedefault
      \SetMathAlphabet \mathit { normal } \g_fontspec_encoding_tl \g__fontspec_mathrm_tl { \mdseries@rm } \itdefault
      \SetMathAlphabet \mathbf { normal } \g_fontspec_encoding_tl \g__fontspec_mathrm_tl { \bfseries@rm } \shapedefault
    }
  \cs_set:Nn \__fontspec_setboldmathrm_hook:nn
    {
      \SetMathAlphabet \mathrm { bold } \g_fontspec_encoding_tl \g__fontspec_bfmathrm_tl { \mdseries@rm } \shapedefault
      \SetMathAlphabet \mathit { bold } \g_fontspec_encoding_tl \g__fontspec_bfmathrm_tl { \mdseries@rm } \itdefault
      \SetMathAlphabet \mathbf { bold } \g_fontspec_encoding_tl \g__fontspec_bfmathrm_tl { \bfseries@rm } \shapedefault
    }
  \cs_set:Nn \__fontspec_setmathsf_hook:nn
    {
      \SetMathAlphabet \mathsf { normal } \g_fontspec_encoding_tl \g__fontspec_mathsf_tl { \mdseries@rm } \shapedefault
      \SetMathAlphabet \mathsf { bold } \g_fontspec_encoding_tl \g__fontspec_mathsf_tl { \bfseries@rm } \shapedefault
    }
  \cs_set:Nn \__fontspec_setmathtt_hook:nn
    {
      \SetMathAlphabet \mathtt { normal } \g_fontspec_encoding_tl \g__fontspec_mathtt_tl { \mdseries@rm } \shapedefault
      \SetMathAlphabet \mathtt { bold } \g_fontspec_encoding_tl \g__fontspec_mathtt_tl { \bfseries@rm } \shapedefault
    }
  %
  \__fontspec_setmathrm_hook:nn {} {}
  \__fontspec_setmathsf_hook:nn {} {}
  \__fontspec_setmathtt_hook:nn {} {}
}

\cs_set_protected:Npn \operator@font {
  \@fontswitch { \font@warning{Math~mode~required~for~\string\operator@font.} } { \mathtextrm }
}

\DeclareSymbolFont {lummain} {TU} {lmm} {m} {n}
\SetSymbolFont {lummain} {bold} {TU} {lmm} {b} {n}

\newattribute \mathfamattr

\lua_load_module:n { lua-unicode-math }

\prop_set_from_keyval:Nn \l_tmpa_prop {
  up = 0, bfup = 1, it = 2, bfit = 3,
  sfup = 4, bfsfup = 5, sfit = 6, bfsfit = 7,
  cal = 8, bfcal = 9,
  frak = 12, bffrak = 13,
  tt = 16,
  bb = 20,
  bf = 1024,
  normal = -"7FFFFFFF,
  literal = -1,
}
\cs_set_eq:NN \mathup \mathrm
\prop_map_inline:Nn \l_tmpa_prop {
  \cs_new_protected:cpn { sym #1 } ##1 {
    {
      \mathfamattr = #2 \scan_stop:
      ##1
    }
  }
  \cs_if_exist:cTF { math #1 } {
    \cs_set_eq:cc { mathtext #1 } { math #1 }
  } {
    \cs_set_eq:cc { math #1 } { sym #1 }
  }
}
\cs_set_eq:NN \mathtextrm \mathrm
\cs_set_eq:NN \symrm \symup
%    \end{macrocode}
% In \pkg{unicode-math} there is a package option to use italic or upright letters
% for \cs{symsf} and \cs{symbfsf}. Depending on the setting these become aliases for
% \cs{symsfup} and \cs{symbfsfup} or \cs{symsfit} and \cs{symbfsfit}.
% Since we do not want package option, we always use upright letters and you ca redefine
% \cs{symsf} and \cs{symbfsf} if you want to change it.
%    \begin{macrocode}
\cs_set_eq:NN \mathtextsf \mathsf
\cs_set_eq:NN \symsf \symsfup
\cs_set_eq:NN \symbfsf \symbfsfup

\clist_map_inline:nn { cal, calbf, frak, frakbf, bb } {
  \cs_set_eq:cc { math #1 } { sym #1 }
}

\cs_new:cpn { __l_uni_math_UnicodeMathSymbol_ \token_to_str:N \mathord :nn } #1 #2 {
  \cs_set:Npx #1 {
    \char_generate:nn {#2} {12}
  }
}
\tl_map_inline:nn {\mathbin \mathclose \mathpunct \mathrel} {
  \cs_new_eq:cc
    { __l_uni_math_UnicodeMathSymbol_ \token_to_str:N #1 :nn }
    { __l_uni_math_UnicodeMathSymbol_ \token_to_str:N \mathord :nn }
}

\cs_new:cpn { __l_uni_math_UnicodeMathSymbol_ \token_to_str:N \mathop :nn } #1 #2 {
  \exp_args:Nc \Umathchardef { \cs_to_str:N #1 op } 1~\symlummain #2~
  \cs_set_eq:cN { \cs_to_str:N #1 oplimits } \scan_stop:
  \cs_set:Npx #1 {
    \char_generate:nn {#2} {12}
  }
  \mathcode #2 = "8000~
  \cs_set:cpx { \char_generate:nn {"FFFF} {12} \char_generate:nn {#2} {12} } {
    \use:c { \cs_to_str:N #1 op }
    \use:c { \cs_to_str:N #1 oplimits }
  }
}

\cs_new:cpn { __l_uni_math_UnicodeMathSymbol_ \token_to_str:N \mathopen :nn } #1 #2 {
  \token_if_eq_meaning:NNTF #1 \sqrt {
    \cs_set:Npx \sqrtsign {
      \Uradical \symlummain #2~
    }
    \cs_set:Npx \root ##1 \of {
      \Uroot \symlummain #2~ { ##1 }
    }
  }{
    \cs_set:Npx #1 {
      \char_generate:nn {#2} {12}
    }
  }
}

\group_begin:
\cs_set:Npn \l_tmp_cs:n #1 {
  \group_end:

  \cs_new_protected:Npn \__l_uni_math__check_mup_helper:w ##1 #1 ##2 \q_mark ##3 ##4 \q_stop {
    ##3 {##2}
  }

  \cs_new:cpn { __l_uni_math_UnicodeMathSymbol_ \token_to_str:N \mathalpha :nn } ##1 ##2 {
    \cs_set:Npx ##1 {
      \char_generate:nn {##2} {12}
    }
    \exp_after:wN \__l_uni_math__check_mup_helper:w \token_to_str:N ##1 \q_mark \cs_set_eq:cN #1 \q_mark \use_none:nn \q_stop ##1
  }
}

\exp_args:No \l_tmp_cs:n {
  \token_to_str:N \mup
}

\cs_new:cpn { __l_uni_math_UnicodeMathSymbol_ \token_to_str:N \mathfence :nn } #1 #2 {
  \cs_set:Npx #1 {
    \char_generate:nn {#2} {12}
  }
  \cs_set:cpx {l \cs_to_str:N #1} {
    \Udelimiter 4 ~ \symlummain #2 ~
  }
  \cs_set:cpx {r \cs_to_str:N #1} {
    \Udelimiter 5 ~ \symlummain #2 ~
  }
}

\cs_new:cpn { __l_uni_math_UnicodeMathSymbol_ \token_to_str:N \mathaccent :nn } #1 #2 {
  \cs_set_protected:Npx #1 {
    \Umathaccent fixed 0 ~ \symlummain #2 ~
  }
}

\cs_new:cpn { __l_uni_math_UnicodeMathSymbol_ \token_to_str:N \mathbotaccent :nn } #1 #2 {
  \cs_set:Npx #1 {
    \exp_not:N \PackageError{lua-unicode-math}{Unsupported~type~\token_to_str:N \mathbotaccent}{}
  }
}

\cs_new:cpn { __l_uni_math_UnicodeMathSymbol_ \token_to_str:N \mathaccentwide :nn } #1 #2 {
  \cs_set:Npx #1 {
    \Umathaccent 0 ~ \symlummain #2 ~
  }
}

\cs_new:cpn { __l_uni_math_UnicodeMathSymbol_ \token_to_str:N \mathbotaccentwide :nn } #1 #2 {
  \cs_set:Npx #1 {
    \Umathaccent bottom 0 ~ \symlummain #2 ~
  }
}

\cs_new:cpn { __l_uni_math_UnicodeMathSymbol_ \token_to_str:N \mathaccentoverlay :nn } #1 #2 {
  \cs_set:Npx #1 {
    \Umathaccent overlay 0 ~ \symlummain #2 ~
  }
  % \cs_set:Npx #1 {
  %   \exp_not:N \PackageError{lua-unicode-math}{Unsupported~type~\token_to_str:N \mathaccentoverlay}{}
  % }
}

\cs_new:cpn { __l_uni_math_UnicodeMathSymbol_ \token_to_str:N \mathover :nn } #1 #2 {
  \cs_set:Npx #1 {
    \exp_not:N \PackageError{lua-unicode-math}{Unsupported~type~\token_to_str:N \mathover}{}
  }
}

\cs_new:cpn { __l_uni_math_UnicodeMathSymbol_ \token_to_str:N \mathunder :nn } #1 #2 {
  \cs_set:Npx #1 {
    \exp_not:N \PackageError{lua-unicode-math}{Unsupported~type~\token_to_str:N \mathunder}{}
  }
}

\cs_generate_variant:Nn \exp_args:Ne {c}
\cs_new:Npn \UnicodeMathSymbol #1 #2 #3 #4 {
  \use:c { __l_uni_math_UnicodeMathSymbol_ \token_to_str:N #3 :nn }
    {#2} {#1}
}
\input {unicode-math-table}
\cs_undefine:N \UnicodeMathSymbol

\cs_set_protected:Npn \triangle { \mathord { \bigtriangleup } }
\cs_set_protected:Npn \mathellipsis { \mathinner { \unicodeellipsis } }
\cs_set_protected:Npn \cdots { \mathinner { \unicodecdots } }

\clist_map_inline:nn {
  \to \rightarrow,
  \le \leq,
  \ge \geq,
  \neq \ne,
  \bigcirc \mdlgwhtcircle,
  \circ \vysmwhtcircle,
  \bullet \smblkcircle,
  \mathyen \yen,
  \mathsterling \sterling,
  \diamond \smwhtdiamond,
  \emptyset \varnothing,
  \hbar \hslash,
  \land \wedge,
  \lor \vee,
  \owns \ni,
  \gets \leftarrow,
  \mathring \ocirc,
  \lnot \neg,
  \longdivision \longdivisionsign,
  \backepsilon \upbackepsilon,
  \eth \matheth,
  \dotsb@ \cdots,
  \@cdots \cdots,
} {
  \cs_set_eq:NN #1
}

\cs_set_eq:NN \intoplimits \nolimits

\cs_set_protected:cpx { \char_generate:nn {"FFFF} {12} ' } {
  \prime_helper:w "2032~
}

\cs_set_protected:Npn \uproot #1 {
  \__l_uni_math_uproot:w #1 \scan_stop:
}

\cs_set_protected:Npn \leftroot #1 {
  \__l_uni_math_leftroot:w #1 \scan_stop:
}
%    \end{macrocode}
% Some fixes for \pkg{amsmath}: Since \pkg{amsmath} is defining \cs{leftroot},
% \cs{uproot} and \cs{root} with non Unicode definitions, we need to hide our definitions
% and restore them afterwards.
% We define \cs{varGamma} to stop \pkg{amsmath} from trying to define greek letter variants.
%    \begin{macrocode}
\tl_const:Nn \c__l_uni_math_amsmath_cmds_tl {
  \uproot
  \leftroot
  \iint
  \iiint
  \iiiint
  \dddot
  \ddddot
  \overleftrightarrow
  \underrightarrow
  \underleftarrow
  \underleftrightarrow
  \hat
  \check
  \tilde
  \acute
  \grave
  \dot
  \ddot
  \breve
  \bar
  \vec
  \mathring
}
\tl_const:Nn \c__l_uni_math_amsmath_cmds_defined_tl {
  \root
  \int
  \oint
  \overrightarrow
  \overleftarrow
}
\hook_gput_code:nnn { package/amsmath/before } {.} {
  \tl_map_inline:Nn \c__l_uni_math_amsmath_cmds_tl {
    \cs_new_eq:cN { __l_uni_math_saved_ \cs_to_str:N #1 } #1
    \cs_undefine:N #1
  }
  \tl_map_inline:Nn \c__l_uni_math_amsmath_cmds_defined_tl {
    \cs_new_eq:cN { __l_uni_math_saved_ \cs_to_str:N #1 } #1
  }
  \cs_set:Npn \varGamma { \temporary_definition_do_not_use }
}
\hook_gput_code:nnn { package/amsmath/after } {.} {
  \tl_map_inline:Nn \c__l_uni_math_amsmath_cmds_tl {
    \cs_set_eq:Nc #1 { __l_uni_math_saved_ \cs_to_str:N #1 }
    \cs_undefine:c { __l_uni_math_saved_ \cs_to_str:N #1 }
  }
  \tl_map_inline:Nn \c__l_uni_math_amsmath_cmds_defined_tl {
    \cs_set_eq:Nc #1 { __l_uni_math_saved_ \cs_to_str:N #1 }
    \cs_undefine:c { __l_uni_math_saved_ \cs_to_str:N #1 }
  }
  \cs_undefine:N \varGamma
}
%    \end{macrocode}
% \end{implementation}
